% !TEX root = ../paper_zh.tex

\section{引言}

红外-可见光对齐是无人机巡检、夜间监测与多模态热感知中的基础能力。可见光影像提供语义与
结构细节，而红外影像可在光照变化、烟雾、薄霾和低照度条件下保留互补的目标信息。两种模态
的结合已经支撑了多光谱检测、跟踪以及复杂运行环境中的航空感知
\cite{hwang2015kaist,flir2018thermal,jia2021llvip,liu2022tardal,sun2022dronevehicle,zhang2022vtuav,li2022lasher,jiang2023antiuav}。
因此，要将两路数据视为一个统一的传感系统，而非彼此独立的观测，就必须准确恢复跨传感器几何关系。

无人机采集引出了一个具有自身特点的对齐问题。跨视场光学系统、不一致的传感器分辨率、变化的
视角以及变焦会使同一场景在可见光和红外通道中呈现不同几何形态。热成像渲染方式以及昼夜外观
变化还会进一步改变局部纹理和对比度。因此，一段飞行序列会提供异质的几何证据：不同图像对在
空间支撑、单应性稳定性以及跨模态一致性方面可能存在显著差异。有效的对齐方法需要将这些证据
整合为可服务于下游感知任务的场景变换。

现代对应关系估计方法已经显著提高了成对证据的可用性。LoFTR、LightGlue、DKM、RoMa、XoFTR
和 MINIMA 等 Transformer、稠密与多模态匹配器，能够在越来越困难的外观变化下恢复对应关系
\cite{sarlin2020superglue,sun2021loftr,chen2022aspanformer,lindenberger2023lightglue,edstedt2023dkm,edstedt2024roma,tuzcuoglu2024xoftr,ren2025minima}。
近期无人机对齐基准也引入了合成变换监督与直接几何评测 \cite{gao2025uavmatch}。这些进展建立了
强大的成对匹配基础，同时也提出了更高层次的基准问题：应如何在真实无人机场景尺度上整合并评估
质量不一的逐帧估计？

本文通过 UAV-TAlign-12K 回答这一问题。该基准包含来自 15 个工业与自然场景的 6,039 对真实采集
跨视场无人机红外-可见光图像。正式评测划分包含 6,037 对经过完整性检查的图像，覆盖白天、夜间
与低照度采集，灰度与伪彩红外渲染，以及广角、变焦和常规无人机视图。该基准定义了三个互补的
评测轨道：成对评测衡量单应性证据可用性；场景级评测衡量多帧共识产品；可靠性--覆盖率曲线刻画
可选择的运行区间。这个层级结构将对应关系生成与无人机感知系统所需的场景产品联系起来。

UAV-TAlign 是该基准的参考框架。它构建在强成对匹配器之上，并结合渐进式证据调度、几何校准
帧加权、鲁棒场景共识和跨模态质量验证。在正式评测划分上，红外专用方法和现代匹配器均能提供
高可用性的成对单应性。UAV-TAlign 将这些证据转化为全部 15 个场景的变换，并在规范运行点识别
出 9 个高置信场景产品，覆盖 74.4\% 的评测图像对。受控扰动和留帧重采样分别从几何敏感性与
证据稳定性两个方面支撑了所得运行曲线。

本文的主要贡献包括：
\begin{itemize}
\item 提出 UAV-TAlign-12K，一个包含 15 个多样化场景、6,039 对同步图像及 6,037 对完整性检查
评测划分的真实跨视场无人机红外-可见光基准。
\item 构建覆盖成对单应性证据、场景级共识与可靠性--覆盖率分析的多层评测协议，并通过条件感知
诊断、重采样稳定性和受控扰动对其进行验证。
\item 提出 UAV-TAlign 参考框架，通过渐进式证据调度、鲁棒共识和跨模态质量验证，将强成对证据
转化为高置信场景产品。
\end{itemize}
